% ===================================================================
%  اللوحة رقم ٩ — الجبل ومدينة المياه. — الغابة والصيد.
%  Traduction 2026 d'apres texto/io, controlee sur texto/fr, texto/en
%  et texto/es. Consigne : voir texto/ar/10-lawha-01.tex.
% ===================================================================

%%K t09-apar-1 apar
\VUsaut{4.0mm}
\VUtitre{15.0pt}{60}{اللوحة رقم ٩}
\VUsaut{3.0mm}

%%K t09-tit-1 sub
\VUcentre{12.0pt}{0}{\textit{المشهد الأول.}}
\VUsaut{3.0mm}

%%K t09-tit-2 sub
\VUpk{11.4pt}{40}{الجبل. --- مدينة المياه.}
\VUsaut{3.0mm}

%%K t09-c1-01-1 p
1. --- أنا في براتز منذ أسبوع. ولا أستطيع أن أصف كل ما أحسست به من عجب
حين وقفت أمام \VUgras{سلسلة الجبال} \textsuperscript{(1)} العجيبة في البرانس،
وقد ارتسم \VUgras{عرفها} \textsuperscript{(2)} على السماء الزرقاء كمنشار عملاق.

%%K t09-c1-02-1 p
2. --- لم أتخيّل أن \VUgras{قمم} \textsuperscript{(3)} البرانس تبلغ هذا
\VUgras{الارتفاع}، وبقيت مذهولًا أمام \VUgras{الأنهار الجليدية}
\textsuperscript{(5)} الرائعة و\VUgras{الذرى} \textsuperscript{(4)} المكسوة بالثلج
التي تنحدر عنها \VUgras{انهيارات} مروّعة.

%%K t09-tit-3 sub
\VUsaut{3.0mm}
\VUpk{10.2pt}{40}{منظر جبلي.}
\VUsaut{3.0mm}

%%K t09-c1-03-1 p
3. --- وفي اليوم التالي لوصولي ذهبت إلى \VUgras{البركان} \textsuperscript{(6)}
القديم الخامد قرب براتز. وعلى حافة \VUgras{الفوهة} العارية الجرداء لا
تزال تُرى بوضوح آثار \VUgras{الحُمَم} التي سالت، منذ قرون، على
\VUgras{سفح الجبل} \textsuperscript{(7)}. وعلى أحد \VUgras{المنحدرات} وُفّقت في
العثور على بضع قطع من تلك الحُمَم.

%%K t09-c1-04-1 p
4. --- وبراتز على الحدود، و\VUgras{القلعة} \textsuperscript{(8)} التي تحرس
\VUgras{الممر} \textsuperscript{(27)} بين فرنسا وإسبانيا تشبه كثيرًا تلك القصور
القديمة على ضفاف الراين التي تبدو كأنها تترصّد \VUgras{الفريسة} من أعلى
صخرتها.

%%K t09-c1-05-1 p
5. --- و\VUgras{نسر} \textsuperscript{(9)} ضخم، وكره غير بعيد عن
\VUgras{الحصن} \textsuperscript{(8)}، كثيرًا ما يشدّ انتباهي. وهذا
\VUgras{الطائر الجارح}، بـ\VUgras{منقاره} القوي، كثيرًا ما يجلب لفراخه
حملًا أو جديًا ممسوكًا في \VUgras{مخالبه}. و\VUgras{جناحاه} من السعة بحيث
يستطيع أن ينساب طويلًا بحمله الثقيل.

%%K t09-tit-4 sub
\VUsaut{3.0mm}
\VUpk{10.2pt}{40}{الماشي.}
\VUsaut{3.0mm}

%%K t09-c1-06-1 p
6. --- وأمتع نزهة هنا هي الذهاب إلى \VUgras{شلال} \textsuperscript{(10)} كو.
\VUgras{الماء الهابط} يتساقط في دوامة ثم يختفي \VUgras{جدولًا} جميلًا في
\VUgras{غابة التنّوب} \textsuperscript{(11)} غربي المدينة. صعدنا هذه الغابة حتى
\VUgras{محطة الأرصاد} \textsuperscript{(12)}، ومن أعلى \VUgras{المرقب} أحطنا
بـ\VUgras{بانوراما} الناحية كلها. \VUgras{المنظر} رائع، ويبدو أن
\VUgras{الطبيعة} أغدقت على هذه البلاد كل ما لديها من كنوز.

%%K t09-c1-07-1 p
7. --- وللنزول ركبنا \VUgras{القطار المائل} \textsuperscript{(13)} ووقفنا عند
\VUgras{محطة} صغيرة قرب \VUgras{أطلال} \textsuperscript{(14)} \VUgras{قصر إقطاعي}
قديم سكنه أمراء براتز فيما مضى.

%%K t09-c1-08-1 p
8. --- أيها القصر المسكين! ماذا عساه يظن وهو يرى تحته \VUgras{مدينة
المياه} \textsuperscript{(15)} الأنيقة وقد صارت \VUgras{مصحًّا حراريًا} من
المصايف الرائجة؟

%%K t09-c1-08-2 p
\VUgras{المناخ} لطيف جدًا و\VUgras{الماء المعدني} نافع جدًا، حتى إن
\VUgras{دورة العلاج} التي تُؤخذ في \VUgras{المصحّ} \textsuperscript{(17)} متعة
حقيقية.

%%K t09-tit-5 sub
\VUsaut{3.0mm}
\VUpk{10.2pt}{40}{مدينة مياه لطيفة.}
\VUsaut{3.0mm}

%%K t09-c1-09-1 p
9. --- والواقع أن كل شيء صُنع ليكون المُقام سارًّا. فقد بُني \VUgras{جسر
علوي} \textsuperscript{(16)} كبير لحمل سكة حديد؛ و\VUgras{حديقة} \textsuperscript{(18)}
\VUgras{قاعة المياه}، حيث تُقام \VUgras{الحفلات الموسيقية}، مرتّبة ترتيبًا
بديعًا. وقد بُنيت عدة \VUgras{فيلات} \textsuperscript{(19)} رشيقة حول
\VUgras{الكازينو} \textsuperscript{(21)}، و\VUgras{طريق عام} \textsuperscript{(22)}
حسن الصيانة جدًا يجعل التنقل في غاية السهولة.

%%K t09-c1-10-1 p
10. --- و\VUgras{سدّ ترابي} \textsuperscript{(23)} بسيط يفصل \VUgras{الطريق}
\textsuperscript{(20)} عن \VUgras{الوادي} \textsuperscript{(25)} نفسه، وفي قاعه
\VUgras{بحيرة} \textsuperscript{(26)} صغيرة رشيقة تغذّي \VUgras{جدولًا}
\textsuperscript{(24)} صافيًا.

%%K t09-c1-11-1 p
11. --- وعلى الجانب الآخر من الوادي يُستغَل \VUgras{منجم} \textsuperscript{(28)}
حديد. وقد ذهبت مرارًا أرى \VUgras{عمال المنجم} ينزلون \VUgras{البئر}، بل
دخلت \VUgras{النفق} حيث يقضون نهارهم في استخراج \VUgras{الخام} الذي يحوي
\VUgras{المعدن}.

%%K t09-c1-12-1 p
12. --- وقد بُني \VUgras{مصهر} كبير قرب المنجم للانتفاع فورًا بما يُستخرج
منه من ثروات. و\VUgras{الفرن العالي} \textsuperscript{(29)}، الذي يُوقَد
بـ\VUgras{الفحم} الوفير هنا، يتيح الحصول على \VUgras{الحديد الخام} سريعًا
ورخيصًا.

%%K t09-c1-13-1 p
13. --- وغير بعيد عن المنجم يُقطع \VUgras{مَقلع} \textsuperscript{(30)}. وليس
ما يقتلعه \VUgras{القطّاع} العجوز بـ\VUgras{مِعوله} \VUgras{غرانيتًا} ولا
\VUgras{سماقًا} ولا حتى \VUgras{حجرًا رمليًا}، بل \VUgras{حجر بناء} عاديًا
لتشييد البيوت.

%%K t09-c1-14-1 p
14. --- ومن أجمل ما يزيّن هذه البلاد \VUgras{شجرة التنّوب} \textsuperscript{(31)}.
ففي كل مكان --- في قاع \VUgras{وادٍ سحيق} \textsuperscript{(32)}، أو على ضفة
\VUgras{سيل} \textsuperscript{(33)}، أو قرب \VUgras{هوّة} \textsuperscript{(34)} أو
منحدر --- تضيف إبرها الرفيعة نغمتها الخضراء.

%%K t09-tit-6 sub
\VUsaut{3.0mm}
\VUpk{10.2pt}{40}{المشي.}
\VUsaut{3.0mm}

%%K t09-c1-15-1 p
15. --- وأنا أغتنم عطلتي لأقوم بنزهات طويلة. فـ\VUgras{الدروب}
\textsuperscript{(35)} التي تتلوّى صاعدة في الجبل تتيح قطع عدة
\VUgras{كيلومترات}، بل عدة \VUgras{فراسخ} في الغالب، دون الشعور بالمسافة.

%%K t09-c1-15-2 p
و\VUgras{أعمدة الأميال} \textsuperscript{(36)} و\VUgras{لافتات الطريق}
\textsuperscript{(37)} تعلّم الدروب، وكثيرًا ما يُلتقى عليها \VUgras{فتى جبلي}
\textsuperscript{(38)} إلى جانبه \VUgras{جِراب} \textsuperscript{(40)} وهو يحمل
\VUgras{غُرَير جبال} \textsuperscript{(39)}، أو \VUgras{غجريّ} \textsuperscript{(41)}
تتنافر \VUgras{قلنسوته الفرائية} \textsuperscript{(42)} تنافرًا غريبًا مع
\VUgras{عباءته} \textsuperscript{(43)} القديمة الممزقة. وهو يقود
بـ\VUgras{سلسلة} \VUgras{دبًّا} \textsuperscript{(44)} مدرَّبًا.

%%K t09-tit-7 sub
\VUpk{10.2pt}{40}{التسلّق.}
\VUsaut{3.0mm}

%%K t09-c1-16-1 p
16. --- وأكاد أذهب كل يوم مع \VUgras{دليل} \textsuperscript{(48)} لأقوم
بـ\VUgras{تسلّقة}، وأنا فخور جدًا حين أحمل \VUgras{حقيبتي} \textsuperscript{(47)}
على ظهري و\VUgras{عصا الجبل} في يدي، كأني \VUgras{سائح} \textsuperscript{(46)}
حقيقي، فأهاجم \VUgras{الصخور} \textsuperscript{(45)}.

%%K t09-c1-17-1 p
17. --- ومن أعلى الجبل لا يبدو أهل السهل أكبر من النمل؛ و\VUgras{قطيع الغنم}
\textsuperscript{(50)} الراعي في \VUgras{مرعى} \textsuperscript{(49)} يبدو
كلعبة طفل. و\VUgras{الصنوبرة} \textsuperscript{(51)}، بل \VUgras{الكوخ الخشبي}
\textsuperscript{(52)}، بالكاد يزيد على نقطة في الخضرة.

%%K t09-c1-18-1 p
18. --- وفي هذه الجولات كثيرًا ما نلتقي على \VUgras{الدرب} \textsuperscript{(53)}
\VUgras{صياد الشمواه} \textsuperscript{(54)} حاملًا على كتفه ما قتله لتوّه.

%%K t09-c1-19-1 p
19. --- وقد وضعت في مجموعتي النباتية سعفة \VUgras{سرخس} \textsuperscript{(55)}
لافتة جدًا وغصينًا ورديًا جميلًا من \VUgras{الخُلَنج} \textsuperscript{(57)}
قطفته عند فم \VUgras{كهف} \textsuperscript{(56)} صغير في نزهتي الأخيرة.

%%K t09-tit-8 sub
\VUsaut{3.0mm}
\VUcentre{12.0pt}{0}{\textit{المشهد الثاني.}}
\VUsaut{3.0mm}

%%K t09-tit-9 sub
\VUpk{11.4pt}{40}{الغابة. --- الصيد.}
\VUsaut{3.0mm}

%%K t09-c2-01-1 p
1. --- بالأمس قضيت في الغابة يومًا بديعًا. وقد شغلتني كثيرًا الحياة
الدائبة للحيوانات و\VUgras{الحشرات} \textsuperscript{(5)} التي تعيش فيها.

%%K t09-c2-01-2 p
\VUgras{العنكبوت} \textsuperscript{(1)} ينسج \VUgras{شبكته} \textsuperscript{(2)} بين
غصنين ليصطاد \VUgras{الذبابة} \textsuperscript{(3)} العابرة، و\VUgras{الحلزون}
\textsuperscript{(4)} يجرّ نفسه بصدفته، والأيّل يستطلع فريسته، والنملة الدائبة
تكدح قرب \VUgras{قرية النمل} \textsuperscript{(6)} --- كل هذه المخلوقات الصغيرة
تروح وتغدو وتعمل بنشاط عجيب.

%%K t09-c2-02-1 p
2. --- و\VUgras{حيّة} \textsuperscript{(7)} حسبتها أول الأمر \VUgras{أفعى} ولم
تكن إلا \VUgras{حيّة ماء} غير مؤذية، كانت ملتفّة تحت جذع شجرة، تخرج بين
الحين والحين رأسها الصغير الرفيع الذي يبدو دائمًا مستعدًا للعضّ. وأمامي
كان \VUgras{بزّاق} \textsuperscript{(8)} كبير يزحف ثقيلًا بعد أن التهم واحدة من
\VUgras{فراولات البر} \textsuperscript{(11)} الصغيرة اللذيذة التي تكثر هنا.

%%K t09-c2-03-1 p
3. --- وكانت الأرض مفروشة بـ\VUgras{أزهار الغابة}: \VUgras{زهر الربيع}
\textsuperscript{(9)} و\VUgras{العَتَرة} \textsuperscript{(10)} ونباتات أخرى كثيرة لم
أكن أعرفها، كلها متفتحة بين \VUgras{خُصَل العشب} \textsuperscript{(12)}.

%%K t09-c2-04-1 p
4. --- وفي هذه الناحية يكاد \VUgras{الكمأ} يكون شائعًا كـ\VUgras{الفطر}
\textsuperscript{(14)}، وكان \VUgras{خنزير صغير} \textsuperscript{(13)} لأحد الحرّاس
ينبش في نهم ليرى أيجد منه شيئًا.

%%K t09-tit-10 sub
\VUsaut{3.0mm}
\VUpk{10.2pt}{40}{الصيد خِلسة.}
\VUsaut{3.0mm}

%%K t09-c2-05-1 p
5. --- ورأيت \VUgras{آخر} مشهد أطربني كثيرًا. فبفضل أنف \VUgras{كلبه
القصير} \textsuperscript{(15)}، كان \VUgras{حارس الغابة} \textsuperscript{(16)} قد
فاجأ لتوّه \VUgras{صيادًا خِلسة} \textsuperscript{(22)} في اللحظة التي كان
يتأهب فيها لحمل \VUgras{الطرائد} \textsuperscript{(17)} التي قتلها. وإذ آثر أن
يتخلى عن صيده على أن يُقبض عليه، فرّ الرجل، وبقي \VUgras{الأرنب}
\textsuperscript{(18)} و\VUgras{الأيّلة} \textsuperscript{(19)} و\VUgras{الغزال}
\textsuperscript{(20)} و\VUgras{الخنزير البري} \textsuperscript{(21)} على العشب،
لسرور الحارس.

%%K t09-c2-06-1 p
6. --- وتنوّع الأشجار في هذه الغابة مدهش. فـ\VUgras{الزان} \textsuperscript{(23)}
و\VUgras{البتولا} \textsuperscript{(26)}، و\VUgras{الصنوبر} الذي يعطي
\VUgras{الراتنج} \textsuperscript{(27)}، و\VUgras{البلوط} \textsuperscript{(29)} وكثير
غيرها، تشابك أغصانها.

%%K t09-c2-06-2 p
و\VUgras{اللبلاب} \textsuperscript{(24)} المتشبث بجذوع الأشجار الباسقة يكسو
\VUgras{الغابة العالية} \textsuperscript{(25)} خضرة زاهية تمتزج امتزاجًا لطيفًا
بلون \VUgras{الطُّحلب} \textsuperscript{(31)} الأشحب الأنعم، حيث يبحث
\VUgras{نقّار الخشب الأخضر} \textsuperscript{(30)} عن الحشرات.

%%K t09-tit-11 sub
\VUsaut{3.0mm}
\VUpk{10.2pt}{40}{منظر الغابة.}
\VUsaut{3.0mm}

%%K t09-c2-07-1 p
7. --- وقد فتنتني أشجار البلوط العتيقة. فـ\VUgras{جذورها} \textsuperscript{(32)}
الكبيرة الملتوية، وقد خرج نصفها من الأرض، تعطيها مظهر قوة ومنعة رائعًا،
وأنا أعجب كيف يطاوع \VUgras{المالك} \textsuperscript{(35)} نفسه على قطعها
وتسليمها إلى \VUgras{منشار} \textsuperscript{(34)} \VUgras{الحطّاب}
\textsuperscript{(33)}. و\VUgras{الجذوع} \textsuperscript{(36)} المسكينة، مقشورة من
\VUgras{لحائها} \textsuperscript{(37)} أو مشقوقة بـ\VUgras{فأس} \textsuperscript{(38)}
\VUgras{الأجير} \textsuperscript{(39)}، تحزنني.

%%K t09-c2-08-1 p
8. --- وعلى مقربة كان \VUgras{فحّام} \textsuperscript{(41)} عجوز قد كوّم
بـ\VUgras{مِجرفته} \VUgras{كومة} \textsuperscript{(42)} من الفحم، وكانت عجوز
تحمل \VUgras{حزمة} \textsuperscript{(40)} من \VUgras{الحطب اليابس} من أغصان
الأشجار المقطوعة.

%%K t09-tit-12 sub
\VUsaut{3.0mm}
\VUpk{10.2pt}{40}{الصيد.}
\VUsaut{3.0mm}

%%K t09-c2-09-1 p
9. --- وكنت أنوي أن أستريح قليلًا في \VUgras{بيت حارس الغابة}
\textsuperscript{(46)}، لكني لما بلغت \VUgras{البِركة} \textsuperscript{(43)} الصغيرة،
وعليها \VUgras{بجعة} \textsuperscript{(45)} جميلة تسبح، لاحظت أن \VUgras{فيضانًا}
\textsuperscript{(44)} حديثًا حوّل الدرب إلى \VUgras{مستنقع} حقيقي. فاضطررت إلى
الالتفاف، وفي مروري بـ\VUgras{الأرزة} \textsuperscript{(49)} و\VUgras{الحور}
\textsuperscript{(48)} و\VUgras{الدردار} \textsuperscript{(47)} القائمة على طرف
الغابة، رأيت يندفع من \VUgras{أجمة} \textsuperscript{(54)}
\VUgras{أيّل} \textsuperscript{(50)} رائع، تطارده \VUgras{زُمرة كلاب} \textsuperscript{(51)}
لا تكلّ، يحثّها \VUgras{بوق} \textsuperscript{(53)} \VUgras{الصيادين على الخيل}
\textsuperscript{(52)}. وكان الحيوان المسكين قد استنفد قواه تمامًا. فسقط يحتضر
على بضع خطوات مني.

%%K t09-c2-10-1 p
10. --- وخرج ابن الحارس من البيت، وقد جذبته ضجة \VUgras{الصيد}، فعدنا
كلانا إلى الغابة. كان قد رأى \VUgras{عقعقًا} \textsuperscript{(56)} على غصن
بلوطة عتيقة. وظن أن عشّه مخبّأ في \VUgras{الأوراق} \textsuperscript{(58)} وأراد
أن يأخذه.

%%K t09-c2-10-2 p
وخفيفًا كـ\VUgras{السنجاب} \textsuperscript{(61)}، تسلّق من \VUgras{غصن}
\textsuperscript{(55)} إلى غصن، لكنه لم يجد شيئًا ولم يُفلح إلا في إفزاع
\VUgras{غراب} \textsuperscript{(60)} عجوز واقف على \VUgras{فرع} \textsuperscript{(57)}.
\VUsaut{4.0mm}
\VUfilet{22mm}
